\documentclass[11pt]{article}
\usepackage[T1]{fontenc}
\usepackage[utf8]{inputenc}
\usepackage{lmodern,amsmath,amssymb,amsthm,mathtools}
\usepackage[margin=1in]{geometry}
\usepackage{microtype,booktabs,array,longtable,enumitem,xcolor}
\usepackage{fancyhdr}
\usepackage[colorlinks=true,linkcolor=blue!45!black,citecolor=blue!45!black,urlcolor=blue!45!black,breaklinks=true]{hyperref}
\usepackage{listings}
\setlength{\emergencystretch}{2em}
\setcounter{tocdepth}{1}
\allowdisplaybreaks[2]
\numberwithin{equation}{section}
\newtheorem{theorem}{Theorem}[section]
\newtheorem{lemma}[theorem]{Lemma}
\newtheorem{proposition}[theorem]{Proposition}
\newtheorem{corollary}[theorem]{Corollary}
\theoremstyle{definition}\newtheorem{definition}[theorem]{Definition}
\theoremstyle{remark}\newtheorem{remark}[theorem]{Remark}
\DeclareMathOperator{\poly}{poly}
\DeclareMathOperator{\supp}{supp}
\DeclareMathOperator{\rank}{rank}
\newcommand{\F}{\mathbb F}
\newcommand{\Z}{\mathbb Z}
\newcommand{\N}{\mathbb N}
\newcommand{\sat}{\#\mathrm{SAT}}
\newcommand{\kcount}{\#k\text{-}\mathrm{SAT}}
\newcommand{\ceil}[1]{\left\lceil#1\right\rceil}
\newcommand{\floor}[1]{\left\lfloor#1\right\rfloor}
\newcommand{\bigO}{\mathcal O}
\pagestyle{fancy}
\fancyhf{}
\fancyhead[L]{\small Exact counting and clause density}
\fancyhead[R]{\small Research note}
\fancyfoot[C]{\thepage}
\setlength{\headheight}{14pt}
\lstset{basicstyle=\ttfamily\small,columns=fullflexible,breaklines=true,frame=single,rulecolor=\color{black!20},xleftmargin=0.5em,xrightmargin=0.5em}
\hypersetup{pdftitle={Improving the Clause-Density Dependence of Exact k-SAT Counting},pdfauthor={ChatGPT research assistance for Vladimir Reshetnikov},pdfsubject={A density-dependent refinement of the October 2026 algorithm}}
\title{Improving the Clause-Density Dependence\protect\\of Exact $k$-SAT Counting}
\author{Research note prepared for Vladimir Reshetnikov\smallskip\\\small ChatGPT research assistance}
\date{October 26, 2026}
\begin{document}
\maketitle
\begin{abstract}
We propose a modification of the exact counting algorithm in OpenAI's
\emph{The Exponential Time Hypothesis is False}, revised October 25, 2026.
Using its Boolean tensor-contraction theorem and parity-lift identity,
and retaining its general counter at high density, we derive, for
fixed $k\ge2$, a deterministic polynomial-space algorithm
whose bit complexity on $n\ge20$ declared variables and $m$ clauses is
\[
 (n+m)^{O_k(1)}
 2^{\,O_k\!\left(\frac{n}{\sqrt{\log n}}
       \min\left\{\log(2+m/n),\log\log n\right\}\right)}.
\]
The change is to construct the initial sparse parity checks using a
powered constant-degree expander with local-code parameters depending
on clause density, and to keep the actual monomial sparsity in the
source's incidence-encoding bound. For $m=\Theta(n\log\log n)$ this
replaces an exponent $O_k(n\log\log n/\sqrt{\log n})$ by
$O_k(n\log\log\log n/\sqrt{\log n})$.
This is an improvement of the density dependence. The unrestricted
worst-case bound in $n$, and the fixed-density exponent, remain those
of the source. We give explicit integer parameters, correctness and
complexity proofs, a boundary analysis, and reproducible finite checks.
The extension is a mathematical derivation relative to the cited
source theorems; it has not been formally verified in Lean or Rocq.
\end{abstract}
\tableofcontents
\clearpage

\section{Result, comparison, and provenance}
\label{sec:result}

All logarithms in this note are base two unless explicitly written
$\ln$. Changing the base affects only fixed constants. The parameter
$k$ is fixed; the density may grow with $n$. The count includes every
declared variable, whether or not it occurs in a clause.

The revised source \cite{OAI2026} states the general bound
\begin{equation}\label{eq:source-general}
 (n+m)^{O_k(1)}2^{O_k(n\log\log n/\sqrt{\log n})}
\end{equation}
and, for each fixed $D$, the bound
\begin{equation}\label{eq:source-sparse}
 (n+m)^{O_{k,D}(1)}2^{O_{k,D}(n/\sqrt{\log n})}
 \qquad(m\le Dn).
\end{equation}
The stronger tensor-contraction appendix, rather than the earlier
rational-interpolation construction, is the baseline here. In fact its
direct clause-network proof also gives
$2^{O_k((n+m)/\sqrt{\log(n+m)})}$ up to a polynomial input factor.
Combining this with \eqref{eq:source-general}, for
$B=\max\{1,\ceil{m/n}\}\le\ceil{\log n}$, gives the exponent scale
\begin{equation}\label{eq:old-profile}
 \frac{n}{\sqrt{\log n}}\min\{1+B,\log\log n\}.
\end{equation}
Keeping the density information through the parity reduction improves
this profile.

\begin{theorem}[Density-dependent improvement]\label{thm:main}
Use the Boolean tensor-contraction theorem, parity-lift identity, and
general-counting bound of \cite{OAI2026} as imported results.
For each fixed integer $k\ge2$, there is a
deterministic algorithm that exactly counts the satisfying assignments
of a $k$-CNF on $n\ge20$ declared variables with $m$ input clauses in
\begin{equation}\label{eq:main}
 (n+m)^{O_k(1)}
 2^{O_k\!\left(\frac{n}{\sqrt{\log n}}
       \min\left\{\log(2+m/n),\log\log n\right\}\right)}
\end{equation}
bit operations and $(n+m)^{O_k(1)}$ bits of space.
All constants are independent of $m/n$.
\end{theorem}

The proof constructs a new reduction for
$B\le\ceil{\log n}$ and uses the source's general algorithm for larger
$B$. A more general version works for $B\le(\log n)^A$ for any fixed
$A$, with constants allowed to depend on $A$. Only $A=1$ is needed
for Theorem~\ref{thm:main}.

\begin{corollary}\label{cor:intermediate}
For $m=\Theta(n\log\log n)$, the improved bound is
\[
 (n+m)^{O_k(1)}
 2^{O_k(n\log\log\log n/\sqrt{\log n})}.
\]
More generally, if $B\to\infty$ and
$\log B=o(\log\log n)$, the new exponent scale is little-oh of
\eqref{eq:old-profile}.
\end{corollary}
\begin{proof}
Substitute the indicated density into \eqref{eq:main}. For the second
claim, both $\log(B+2)/(1+B)$ and
$\log(B+2)/\log\log n$ tend to zero. The ratio to the minimum
in \eqref{eq:old-profile} is the maximum of these two ratios.
\end{proof}

\paragraph{What is inherited and what changes.}
The contraction algorithm and its $s/\sqrt{\log s}$ exponent are
inherited. The parity-lift identity and the monomial-to-network
translation are also inherited. The new ingredients in this comparison
are the density-only bounds on check support, and the uniform retention
of the actual sparsity in the network-size estimate. Expander codes
themselves are classical and already appear in the source. This note
does not claim a new coding paradigm or priority over unexamined work.

\paragraph{Version examined.}
The source snapshot is \texttt{openai/math} at commit
\begin{center}
\small\texttt{7ef5c0cc5ca4b0e072b7a050fca09ab24bc6fa3a}.
\end{center}
The analysis uses three sources in the preprint's \texttt{build/} directory:
\begin{itemize}[nosep,leftmargin=*]
\item \texttt{main\_full\_version.tex};
\item \texttt{quantitative-reduction.tex};
\item \texttt{quantitative-contraction.tex}.
\end{itemize}
The accompanying source manifest records
their Git blob hashes and permanent URLs. This note does not report
rerunning the source's formal verification or implementing its complete
contraction engine.

\section{A sparse-polynomial counter that retains its sparsity}
\label{sec:sparse}

The first step is a quantitative refinement of the source's lemma
headed ``Counting sparse polynomial systems by contraction.'' Its
proof gives more information than its stated polylogarithmic-sparsity
bound when the actual number of monomials is retained.

The imported contraction interface allows a bipartite Boolean tensor
network with $O(s)$ edges and maximum degree at most $\ceil{s^{9/10}}$,
over $(\Z/2^p\Z)[V]/(V^\ell)$ with $p,\ell=O(s)$. There must be
polynomially many vertices and a description of length $\mathcal L$.
If each tensor entry is computable in
$(s+\mathcal L)^{O(1)}2^{O(s^{9/10})}$ time and
$(s+\mathcal L)^{O(1)}$ space, the source theorem contracts the network
in $(s+\mathcal L)^{O(1)}2^{O(s/\sqrt{\log s})}$ time and polynomial
space. All constants in this interface are fixed. This is the input
whose hypotheses are checked below.

\begin{lemma}[Uniform sparsity dependence]\label{lem:sparse}
Fix positive integers $k,c,C$. Let $P_1,\ldots,P_a$ be polynomials
over $\F_2$ on $n\ge20$ declared Boolean variables, where $a\le cn$.
Each polynomial is given by at most $S$ monomials of degree at most
$k$, with
\[
 1\le S\le\lambda^C,\qquad \lambda=\ceil{\log n}.
\]
Under the source assumptions of Theorem~\ref{thm:main}, their number
of common zeros can be computed deterministically in polynomial space
and
\begin{equation}\label{eq:sparse-time}
 n^{O_{k,c,C}(1)}
 2^{O_{k,c,C}(n\log(S+2)/\sqrt{\log n})}
\end{equation}
bit operations. The constants are uniform over all permitted $S$.
\end{lemma}

\begin{proof}
Replace positive powers of each Boolean variable by its first power,
and cancel duplicate monomials modulo two. There are at most
\begin{equation}\label{eq:incidences}
 N_0\le akS\le cknS
\end{equation}
equation--variable incidences. The scope of one equation has size at
most $kS$.

\paragraph{Control the large variable degrees.}
Set
\[
 D_* = \max\{1,\ceil{N_0/n}\}\ceil{\sqrt\lambda}.
\]
Branch over all assignments to variables belonging to more than
$D_*$ scopes. There are at most $n/\sqrt\lambda$ such variables, by
counting incidences. Thus branching contributes at most
$2^{n/\sqrt\lambda}$ calls. Substitute the fixed bits, normalize
again, and discard constant-zero equations. A constant-one equation
gives zero immediately. All remaining variable incidences are bounded
by $O_{k,c}(S\sqrt\lambda)$; every remaining equation still has at
most $S$ monomials. Keep the original value of $n$ for all bounds and
precisions.

\paragraph{Use the source's bounded-degree lift.}
The source's parity-lift construction works over
\[
 \mathcal B=(\Z/2^{n+1}\Z)[V]/(V^{n+1}).
\]
Give each surviving incidence $(i,v)$ its own variable $z_{iv}$.
The construction produces a multilinear polynomial
\[
 F(z,V)=\prod_i \overline B_i(z_i,V)
\]
with two properties. First, every retained incidence monomial has
degree at most
\begin{equation}\label{eq:degree}
 K=k(a+2n),\qquad 2kn\le K\le(c+2)kn.
\end{equation}
Second, after identifying every $z_{iv}$ with the original Boolean
$x_v$, summing over the free Boolean assignments, and summing the
$V$-coefficients, the result is congruent modulo $2^{n+1}$ to the
branch count. These are precisely the degree and residue identities
in the source's sparse-query proof. They use its total-order parity
lift; no new parity approximation is being assumed here.

\paragraph{Encode each monomial by subsets of incidences.}
Partition the tests and the free variables separately into
$J=\ceil{n^{1/4}}$ groups, balanced in cardinality. Let $N_{gh}$
be the number of incidences between test group $g$ and variable group
$h$, and let $N=\sum_{g,h}N_{gh}\le N_0$. Enumerate arrays
$(d_{gh})$ with $0\le d_{gh}\le N_{gh}$ and
$\sum d_{gh}\le K$. There are at most
\[
 (K+1)^{J^2}=2^{O_{k,c}(\sqrt n\log n)}
\]
such arrays. For a fixed array, use
$\ceil{\log\binom{N_{gh}}{d_{gh}}}$ Boolean wires to encode the
chosen incidence subset for each pair of groups. Invalid indices
receive tensor value zero. The source's tensor entries multiply the
indicated test coefficients and sum the identified variable powers;
therefore summing these networks over the degree arrays computes the
required lift sum exactly.

The important estimate is the number of wires $M$:
\begin{align}
 M&\le J^2+\log\prod_{g,h}\binom{N_{gh}}{d_{gh}}\notag\\
  &\le J^2+\log\binom{N+K}{K}\notag\\
  &\le J^2+K\log\!\left(\frac{\mathrm e(N+K)}K\right)
    =O_{k,c}\bigl(n\log(S+2)\bigr).
 \label{eq:wire-entropy}
\end{align}
The second line bounds the number of the specified squarefree
monomials by the number of all monomials of degree at most $K$ in
$N$ variables. The last line uses \eqref{eq:incidences} and
\eqref{eq:degree}. Replacing $S$ by an unspecified power of $\log n$
at this point would discard the improvement.

\paragraph{Check the contraction hypotheses uniformly.}
Choose the integer size parameter
\[
 \tau=1+\ceil{\log(S+1)},\qquad s=n\tau.
\]
Then \eqref{eq:wire-entropy} gives $M\le C_0s$ for a fixed
$C_0=C_0(k,c)$. There are $2J$ vertices. Their degrees are at most
\[
 O_{k,c}(n^{3/4}S\sqrt\lambda)=o(s^{9/10}),
\]
uniformly for $S\le\lambda^C$. Each test-group tensor entry is
computed by subset M\"obius inversion, one test at a time, in
$2^{kS}n^{O_{k,c,C}(1)}$ time and polynomial space. This is at most
$n^{O_{k,c,C}(1)}2^{s^{9/10}}$ for all sufficiently large $n$,
uniformly over the same range. Variable-group entries and index
decoding have polynomial cost. The graph and the entry descriptions
have polynomial length; the precision and truncation order are
$n+1=O(s)$.

The source tensor theorem thus applies with fixed constants. It
gives time polynomial in the description length times
\[
 2^{O_{k,c,C}(s/\sqrt{\log s})}
  =2^{O_{k,c,C}(n\log(S+2)/\sqrt{\log n})}.
\]
Both the high-incidence branching exponent $n/\sqrt\lambda$ and
the array-enumeration exponent $O(\sqrt n\log n)$ are absorbed in
this bound. Process branches, arrays, and contractions sequentially.
The source's depth-first polynomial-space argument applies unchanged.
The final count lies in $[0,2^n]$, so its least nonnegative residue
modulo $2^{n+1}$ is the exact integer count. A finite range of smaller
$n$, depending only on $k,c,C$, can be handled by direct enumeration.
This changes only fixed constants.
\end{proof}

\section{Initial parity checks with support depending only on density}
\label{sec:checks}

The source's initial checks use, with $m\le Bn$ and
$\mu=\floor{\log m}+1$,
\[
 D_0=320000B^2\mu^3,\qquad P_0=20000B^2\mu^2.
\]
Here $D_0$ bounds initial row supports and $P_0$ is the later sampling
length. These factors of $\mu$ make the query sparsity a power of
$\log n$ even when $B$ grows very slowly. We replace this part of
the construction.

\subsection{The explicit expander input}

For every integer $q\ge1$, the Margulis--Gabber--Galil graph on
$(\Z/q\Z)^2$ has the following eight labelled ports from $(x,y)$:
\begin{equation}\label{eq:mgg}
 (x\pm2y,y),\quad(x\pm(2y+1),y),\quad
 (x,y\pm2x),\quad(x,y\pm(2x+1)).
\end{equation}
All operations are modulo $q$. Parallel edges and self-loops retain
their port identities. The classical spectral theorem bounds the
second largest eigenvalue of its normalized adjacency matrix $A/8$
by $5\sqrt2/8$ \cite{GG1981,MGGFormal}. The cited formalization
also records the explicit maps and the spectral proof.

Add eight labelled identity ports and use the degree-$16$ lazy
operator
\[
 Q=\frac{I+A/8}{2}.
\]
Every eigenvalue is nonnegative, and its restriction to the
orthogonal complement of the constants has norm at most
\begin{equation}\label{eq:rho}
 \rho=\frac{1+5\sqrt2/8}{2}<\frac{19}{20}.
\end{equation}
This laziness avoids any convention about whether ``second
eigenvalue'' bounds an absolute value. In the $t$-step walk graph,
the degree is $16^t$ and the corresponding norm is at most
$(19/20)^t$. A port is a word of length $t$ in the $16$-symbol
alphabet. Each step map is a permutation, so there are exactly
$16^t$ incoming and outgoing walk ports at every vertex.

\subsection{Integer parameters and the small-instance branch}

Let $B\ge1$ be an integer with $m\le Bn$. Define
\begin{equation}\label{eq:parameters}
 b=\ceil{\log(2B)},\qquad A_B=4096Bb,\qquad
 \delta=A_B^{-1},\qquad P=2A_B^2.
\end{equation}
Choose the least integer $t\ge1$ satisfying
\begin{equation}\label{eq:integer-tests}
 2A_B19^t\le20^t,\qquad16^t\ge64B,
 \qquad\text{and set }d=16^t.
\end{equation}
No real-number oracle is needed. These integer inequalities imply
$\varepsilon=(19/20)^t\le\delta/2$ and $d\ge64B$.
They also give
\begin{equation}\label{eq:d-bound}
 t=O(\log(B+2)),\qquad d\le(16384Bb)^{56}.
\end{equation}
Indeed $2\cdot19^{14}<20^{14}$, so
$t=14\ceil{\log(8192Bb)}$ is an admissible, possibly nonminimal
choice. Since $2^{\ceil{\log x}}<2x$ for $x>0$ except at harmless
equality endpoints, its degree is bounded as in \eqref{eq:d-bound}.
This same choice also satisfies $16^t\ge64B$.

If $m\le d$, directly enumerate the assignments of the at most
$km$ occurring variables and multiply the resulting count by
$2^{n-n_{\rm occ}}$. This costs
\begin{equation}\label{eq:small-m}
 \poly_k(n,m)2^{kd}=\poly_k(n,m)2^{(B+2)^{O_k(1)}}.
\end{equation}
When $B\le\ceil{\log n}$ this is an asymptotically subordinate
overhead. It also avoids padding a tiny formula to an enormous graph.
The constants in \eqref{eq:d-bound} are deliberately conservative;
they are not intended to provide a competitive practical solver.

Assume for the rest of this section that $m>d$. Set
\begin{equation}\label{eq:padding}
 q=\ceil{\sqrt{m/d}},\qquad N=q^2,\qquad M=Nd.
\end{equation}
The integer $q$ can be found as the least positive integer satisfying
$q^2d\ge m$. Since $m/d>1$, we have
\begin{equation}\label{eq:padding-bound}
 m\le M<4m.
\end{equation}
Take a left and a right copy of the $N$ vertices of the powered lazy
graph. Each labelled walk gives an edge from its left starting
vertex to its right ending vertex. This is a $d$-regular bipartite
multigraph. Spectral decomposition gives, for left and right vertex
sets $U,V$,
\begin{equation}\label{eq:mixing}
 e(U,V)\le \frac{d|U||V|}{N}
              +\varepsilon d\sqrt{|U||V|}.
\end{equation}
To obtain it, subtract the constant components of the two indicator
vectors and bound the residual bilinear form by its operator norm.
Replacing the residual norms by $\sqrt{|U|}$ and $\sqrt{|V|}$
only weakens the estimate.

\subsection{A local code with density-dependent distance}

\begin{lemma}\label{lem:local-code}
There exists a binary matrix $H$ with $d$ columns and
$c_{\rm loc}=\ceil{d/(64B)}$ rows such that every nonzero word in
its kernel has Hamming weight at least $\delta d$. Such a matrix
can be found deterministically in $2^{O(d^2)}$ bit operations and
$O(d^2)$ working bits.
\end{lemma}
\begin{proof}
For $0<x\le1/2$, let
$H_2(x)=-x\log x-(1-x)\log(1-x)$ be binary entropy. The usual
binomial bound and $H_2(x)\le x\log(\mathrm e/x)$ give
\[
 \sum_{j=0}^{\ceil{\delta d}-1}\binom dj
      \le2^{dH_2(\delta)}.
\]
Using $\log B\le b-1$, $\log b\le b$, $b\ge1$, and
$\log\mathrm e<2$, we have
\[
 \log(\mathrm e/\delta)
   =\log\mathrm e+12+\log B+\log b
   <13+2b\le16b.
\]
Consequently
\begin{equation}\label{eq:local-entropy}
 H_2(\delta)\le\frac1{256B}
     <\frac1{64B}\le\frac{c_{\rm loc}}d.
\end{equation}
For a uniformly random $c_{\rm loc}$-by-$d$ binary matrix, every
fixed nonzero vector belongs to its kernel with probability
$2^{-c_{\rm loc}}$. The expected number of nonzero kernel words
of weight below $\delta d$ is strictly less than one by
\eqref{eq:local-entropy}. Some matrix has no such word. Enumerate
the matrices lexicographically, testing the forbidden words for each.
The total cost is at most $2^{c_{\rm loc}d+d}\poly(d)$, which is
$2^{O(d^2)}$, and only the current matrix and word are needed in
working storage. The test $w<\delta d$ is the exact integer test
$A_Bw<d$.
\end{proof}

\subsection{Global checks, distance, and shortening}

Place the local checks $H$ on the incident ports of every vertex on
both sides of the bipartite graph. The resulting global matrix has
\begin{align}
 2Nc_{\rm loc}
   &\le\frac{M}{32B}+\frac{2M}{d}\notag\\
   &\le\frac{M}{16B}<\frac{m}{4B}\le\frac n4
 \label{eq:row-count}
\end{align}
rows. Each row reads at most $d$ edge bits. The padding factor in
\eqref{eq:padding-bound} is explicitly paid for in this estimate.

\begin{lemma}\label{lem:initial-checks}
When $m>d$, the construction gives a deterministic binary matrix
$A^{(0)}$ with $m$ columns and fewer than $n/4$ rows, each of
support at most $d$, satisfying
\begin{equation}\label{eq:initial-distance}
 0\ne z\in\ker_{\F_2}A^{(0)}
       \quad\Longrightarrow\quad |\supp z|\ge m/P.
\end{equation}
Here $d,P$ are bounded by fixed polynomials in $B+2$.
\end{lemma}
\begin{proof}
First consider a nonzero accepted word on all $M$ edges. Let $w$
be its weight and $U,V$ the vertex sets touched by its nonzero
edges. The local-code property gives
\[
 w\ge\delta d\max\{|U|,|V|\}
       \ge\delta d\sqrt{|U||V|}.
\]
Every nonzero edge joins $U$ to $V$, so $w\le e(U,V)$.
Combining with \eqref{eq:mixing} and dividing by the positive number
$d\sqrt{|U||V|}$ yields
\[
 \sqrt{|U||V|}\ge(\delta-\varepsilon)N.
\]
It follows that
\[
 w\ge\delta(\delta-\varepsilon)Nd
     \ge\frac{\delta^2M}{2}\ge\frac{m}{P}.
\]
Keep any $m$ labelled edge coordinates and fix the remaining ones
to zero. Restrict every global row to the retained columns to obtain
$A^{(0)}$. A word in the restricted kernel extends by zeros to a
word in the original kernel. Its weight is unchanged, so the distance
lower bound survives this shortening. Row count and support size
cannot increase. Equations \eqref{eq:parameters} and
\eqref{eq:d-bound} prove the claimed polynomial dependence on $B$.
\end{proof}

Constructing the graph and its port lists uses polynomial time in
$n,m,d,t$. Searching for the local matrix costs $2^{O(d^2)}$ once.
The graph has fewer than $4m$ edges. The whole construction uses
space polynomial in $n,m,d$, hence polynomial in $n,m$ in the
density range used below.

\section{From sparse checks to exact counting}
\label{sec:greedy}

For a clause $e$, let $v_e(x)\in\F_2$ indicate its violation by
assignment $x$. As a Boolean polynomial it has degree at most $k$
and at most $2^k$ monomials. Let $v(x)=(v_1(x),\ldots,v_m(x))$.
Given a binary check matrix $A$, define
\[
 N(A)=\bigl|\{x\in\{0,1\}^n:Av(x)=0\}\bigr|.
\]
Every satisfying assignment survives every linear check. Therefore
\begin{equation}\label{eq:survivors}
 N(A)=\#F+D(A),
\end{equation}
where $D(A)$ counts the surviving nonsatisfying assignments.

\begin{lemma}[Greedy elimination]\label{lem:greedy}
Starting with $A^{(0)}$ from Lemma~\ref{lem:initial-checks}, a
deterministic oracle reduction appends $3n$ rows of list length $P$
and obtains $N(A)=\#F$. It makes at most $3n(2m)^P+1$ queries to
an exact common-zero counting oracle. Every query has fewer than $4n$ equations, of
degree at most $k$, and at most
\begin{equation}\label{eq:query-sparsity}
 S_B=2^k\max\{d,P\}
\end{equation}
monomials in each equation.
\end{lemma}
\begin{proof}
Represent a candidate row by an ordered list
\[
 \gamma=((e_1,c_1),\ldots,(e_P,c_P))
       \in([m]\times\{0,1\})^P.
\]
Repeated clause indices are allowed; their coefficients are added
modulo two. Fix a nonsatisfying survivor $x$ of the current matrix.
Its violation vector remains in the kernel of $A^{(0)}$ and has
support at least $m/P$. A uniform candidate list misses this support
with probability at most
\[
 (1-1/P)^P\le e^{-1}<1/2.
\]
Conditional on a hit, an independent coefficient at a hit position
makes the parity uniform, even if indices repeat. Thus at least
one quarter of the candidate lists reject $x$. Summing over all
nonsatisfying survivors gives
\[
 \mathbb E_\gamma D(A;\gamma)\le\tfrac34D(A).
\]
Enumerate every candidate and compute $N(A;\gamma)$ using the
exact common-zero counting oracle. Append a minimizing row. Equation
\eqref{eq:survivors} shows that this also minimizes $D$ even though
$\#F$ is not known. After $3n$ rounds,
\[
 D(A)\le(3/4)^{3n}2^n=(27/32)^n<1.
\]
Since $D(A)$ is a nonnegative integer, it is zero. One final query
returns $\#F$.

The initial row count is below $n/4$; at most $3n$ rows are ever
appended. Every tested system therefore has fewer than $4n$
equations. A row supported on at most $d$ clauses expands to at
most $2^kd$ monomials, and a list of length $P$ expands to at most
$2^kP$ monomials. Cancellation only decreases these numbers.
This proves \eqref{eq:query-sparsity} and the query bound.
\end{proof}

The argument is the source's greedy selection step with new support
parameters. Its use of exact common-zero counts is not circular:
Lemma~\ref{lem:sparse} depends on the already available tensor
counter, not on the density-dependent SAT algorithm being constructed.

\section{The full algorithm and its uniform complexity}
\label{sec:algorithm}

\begin{proof}[Proof of Theorem~\ref{thm:main}]
Normalize the input by deleting tautologies and repeated clauses and
literals, without changing the list of declared variables. An empty
clause gives zero; a formula with no clauses gives $2^n$. Let $m'$
be the number of remaining clauses and
$B=\max\{1,\ceil{m'/n}\}$. Since $m'\le m$, bounds expressed in
terms of $m'$ are no larger, up to the input-processing polynomial,
than the claimed bounds in terms of $m$.

Suppose first that $B\le\lambda=\ceil{\log n}$. Use
\eqref{eq:parameters}--\eqref{eq:integer-tests}. If $m'\le d$, use
the direct occurring-variable branch \eqref{eq:small-m}. Otherwise
construct the checks and use Lemma~\ref{lem:greedy}.
By \eqref{eq:d-bound},
\begin{equation}\label{eq:uniform-s}
 S_B=(B+2)^{O_k(1)},\qquad
 \log(S_B+2)=O_k(\log(B+2)).
\end{equation}
For one fixed exponent $C=C(k)$, all these sparsities are at most
$\lambda^C$ for all sufficiently large $n$. Lemma~\ref{lem:sparse}
with $c=4$ therefore gives query exponent
\begin{equation}\label{eq:new-query-exp}
 O_k\!\left(\frac{n\log(B+2)}{\sqrt{\log n}}\right),
\end{equation}
with constants independent of the actual value of $B$.

The logarithm of the number of queries, the local-code search cost,
and the small-instance enumeration cost are bounded by
\begin{equation}\label{eq:overheads}
 O_k\bigl(P\log(2m')+d^2+kd+\log n\bigr)
      =\lambda^{O_k(1)}.
\end{equation}
For every fixed $C'$, $\lambda^{C'}=o(n/\sqrt\lambda)$.
Hence \eqref{eq:overheads} is absorbed in \eqref{eq:new-query-exp},
uniformly over $B\le\lambda$. Multiplication by the number of
queries adds their logarithm to the exponent; it does not multiply
the exponential scales. The time is therefore
\begin{equation}\label{eq:density-only}
 (n+m)^{O_k(1)}
 2^{O_k(n\log(B+2)/\sqrt{\log n})}.
\end{equation}

If $B>\lambda$, use the source algorithm
\eqref{eq:source-general}. In this range
$\log(2+m'/n)=\Omega(\log\log n)$, so its exponent agrees, up to
a fixed factor, with the minimum in \eqref{eq:main}. In the first
range $\log(B+2)=\Theta(\log(2+m'/n))$ and it is at most a fixed
multiple of $\log\log n$. Thus \eqref{eq:density-only} also has the
form \eqref{eq:main}. These comparisons prove the time bound.

For space, enumerate matrices and candidates one at a time, keep
only the current and best candidate counts, and use the
polynomial-space sparse-system counter. The selected matrix has
$O(n)$ rows; each list or support has polylogarithmic length in the
first range. The graph has fewer than $4m'$ edges; local-code
enumeration stores $O(d^2)$ bits. Accumulators use $n+1$ bits for
counts. The direct occurring-variable branch is depth first, and
the source's general branch already has polynomial space. All
claimed space bounds follow. A uniformly finite exceptional range
of $n$ is handled directly and absorbed in the fixed constants.
\end{proof}

The same low-density construction works, for any fixed $A>0$, on
$B\le\lambda^A$ with constants depending on $k,A$. Every support,
local search, and candidate-enumeration overhead is still bounded by
a fixed power of $\lambda$ in the exponent. This proves the
fixed-$A$ extension mentioned after Theorem~\ref{thm:main}.

\subsection{Algorithm in operational form}

\begin{enumerate}[leftmargin=*,itemsep=0.35em]
\item Normalize the CNF, preserving its $n$ declared variables. Handle
      an empty clause or no remaining clauses immediately.
\item Compute $B=\max\{1,\ceil{m'/n}\}$ and
      $\lambda=\ceil{\log n}$. If $B>\lambda$, call the source's
      general counter and return its result.
\item Compute the exact integer parameters in
      \eqref{eq:parameters}--\eqref{eq:integer-tests}. If $m'\le d$,
      count by enumerating the occurring variables, include the
      unused-variable factor, and return.
\item Construct the powered lazy expander, search for its local
      check matrix, impose local checks at both ends, and shorten
      the edge code to $m'$ coordinates.
\item For $3n$ rounds, enumerate lists in $([m']\times\{0,1\})^P$.
      Expand each candidate into degree-$k$ equations and count
      their common zeros with Lemma~\ref{lem:sparse}. Append a
      candidate giving the minimum count.
\item Count the common zeros of the final system and return that
      integer. Lemma~\ref{lem:greedy} proves it equals $\#F$.
\end{enumerate}

No randomized choice remains in this algorithm. Random matrices and
random lists are used only to prove that exhaustive deterministic
searches find objects with the required properties.

\section{Where the bound improves, and where it does not}
\label{sec:scope}

Table~\ref{tab:comparison} compares the exponential scales after
factoring out $n/\sqrt{\log n}$. Polynomial factors and fixed
constants are suppressed. These are comparisons of proved upper
bounds, not runtime measurements or lower bounds on the source.

\begin{table}[htbp]
\centering
\caption{Clause-density dependence of the exponent.}
\label{tab:comparison}
\begin{tabular}{@{}p{0.30\linewidth}p{0.29\linewidth}p{0.28\linewidth}@{}}
\toprule
Density $B$ & Source guarantees combined & This modification\\
\midrule
$\Theta(1)$ & $1$ & $1$\\[0.4em]
$\sqrt{\log\log n}$ & $\sqrt{\log\log n}$ & $\log\log\log n$\\[0.4em]
$\log\log n$ & $\log\log n$ & $\log\log\log n$\\[0.4em]
$(\log\log n)^2$ & $\log\log n$ & $\log\log\log n$\\[0.4em]
$(\log n)^\alpha$, fixed $\alpha>0$ & $\log\log n$ & $\log\log n$\\[0.4em]
Unrestricted & $\log\log n$ & $\log\log n$\\
\bottomrule
\end{tabular}
\end{table}

For the central example $B=\Theta(\log\log n)$, the ratio of the
old and new exponent scales tends to infinity like
\[
 \frac{\log\log n}{\log\log\log n}.
\]
This is more than a constant-factor change in the exponential
upper bound. It is also narrower than improving the worst case over
all $m$: at unrestricted density, \eqref{eq:main} retains the source
exponent. Nor does it improve the fixed-density
$n/\sqrt{\log n}$ scale.

\subsection{Why simply retuning sparsification does not go further}

The source's disjoint sparsification step, with parameter
$0<\sigma\le1$, incurs $2^{\sigma n}$ leaves up to polynomial
factors, with an occurrence bound
\[
 B=O_k\bigl([\sigma^{-1}\log(2/\sigma)]^{k-1}\bigr).
\]
Substituting this bound into \eqref{eq:density-only} gives the
exponent estimate
\begin{equation}\label{eq:sparsify-objective}
 O_k\!\left(n\left[
       \sigma+\frac{\log(2/\sigma)}{\sqrt{\log n}}
                    \right]\right).
\end{equation}
On its permitted polylogarithmic-density range, the natural minimizing
scale is $\sigma=\Theta(1/\sqrt{\log n})$. For example, minimizing
$\sigma+a\ln(1/\sigma)/\sqrt{\log n}$ for a positive fixed $a$
gives $\sigma=a/\sqrt{\log n}$. Its value has scale
$\log\log n/\sqrt{\log n}$. Thus this particular upper-bound
calculation still has the source's general scale. This observation
is a limitation of the present estimate; it is not an impossibility
theorem for other sparsifications or other algorithms.
For general fixed $k$, this substitution uses the fixed-$A$
extension with $A$ chosen sufficiently large in terms of $k$:
the occurrence bound at the indicated $\sigma$ is polylogarithmic,
although it can exceed $\log n$.

\subsection{Why a total-degree cutoff alone does not beat the tensor exponent}

The source's sharper tensor routine enumerates weighted mode
monomials. With $M$ wires, $u$ retained valuation levels, and a
packet parameter $D$, its relevant upper budgets have scales
\[
 \text{colors per weight }O(M/D),\qquad
 \text{weight budget }O(MD/u),
\]
giving $\exp(O(M/\sqrt u))$ tuples. Here $u=\Theta(\log s)$.
A tempting suggestion is to additionally impose the ordinary
degree bound $|\beta|\le M$. The following elementary example shows
why those two budgets by themselves cannot yield a better
asymptotic enumeration bound.

\begin{proposition}[A counting obstruction, not an algorithmic lower bound]
\label{prop:obstruction}
Let $M,D,u$ be positive integers. Suppose an abstract mode family
has $\floor{M/D}$ distinct modes
of every weight in an interval $[b,2b]$, where
$b=\floor{D/(4\sqrt u)}$, $M/D\ge1$, and $D/\sqrt u$ is
sufficiently large. Then there are $2^{\Omega(M/\sqrt u)}$
distinct exponent tuples satisfying both $|\beta|\le M$ and
$\sum_jw_j\beta_j<MD/u$.
\end{proposition}
\begin{proof}
There are $\Theta((M/D)b)=\Theta(M/\sqrt u)$ modes in this
interval. Give each of them an exponent in $\{0,1\}$ independently
and all other modes exponent zero. Even selecting them all uses at
most $\floor{M/D}(b+1)\le M$ total degree for sufficiently large
parameters. The total weight is at most
\[
 2b\floor{M/D}(b+1)<MD/u
\]
for the chosen constant $1/4$, again for sufficiently large
parameters. Every subset is therefore allowed, yielding the stated
number of tuples.
\end{proof}

This proposition does not assert that all such tuples are nonzero
in every actual network. It shows that the two numerical budgets
alone cannot justify an improved counting estimate. Additional
algebraic structure or a way to sum many modes without enumerating
them would be needed.

\section{Verification status and reproducible artifacts}
\label{sec:verification}

The proof was checked independently at the following interfaces:
the expander padding and port counts; the strict local-code entropy
inequality; the global row budget; the minimum-distance argument
under shortening; the greedy elimination invariant; and the
uniform contraction hypotheses when the density grows with $n$.
The source's complete tensor proof and its formalization were not
reverified in this work.

The archive includes a standard-library-only Python verification
script and its machine-readable results. The script performs exact
finite checks of the integer parameter inequalities, the explicit
graph maps and their inverse ports, small spectral bounds, local and
shortened code examples, subset indexing, and greedy parity
selection. It also records the scale of the explicit parameters.
These checks support implementation details and catch bookkeeping
errors; the unbounded assertions are proved in the text or imported
from the cited theorems.

The executed check suite passed: $3{,}248$ inverse-port identities;
four exact rational spectral certificates for side lengths $2$ through
$5$; five small local-code searches; $156$ shortening and zero-extension
checks; $2{,}047$ subset rank/unrank cases; $30{,}000$ wire-encoding
inequalities; $72$ greedy parity rounds; parameter choices for
$B=1,\ldots,1000$; and $27{,}696$ padding and row-budget cases.
The spectral certificates use exact rational elimination, not
floating-point eigenvalue approximations. These figures describe the
finite tests only.

\paragraph{What the code is.}
It is a reference implementation and verifier for the new reduction's
components. It is not a complete implementation or benchmark of the
subexponential SAT counter. In particular, it does not replace the
source's tensor engine by brute force and then claim the engine's
asymptotic running time for that brute-force code.

\paragraph{Reproduction.}
From the extracted archive directory, run:
\begin{lstlisting}
python3 code/verify_components.py --output verification_results.json
latexmk -pdf -interaction=nonstopmode -halt-on-error density_aware_sat_improvement.tex
\end{lstlisting}
The manuscript uses only common LaTeX packages. The Python checks
require no third-party packages. The source manifest identifies the
upstream files used; \texttt{PROOF\_STATUS.txt} records the proof
dependencies and the distinction between finite tests and formal
verification. \texttt{README.txt} describes each artifact.

\paragraph{Practical limitation.}
Even at small $B$, the conservative choice of $t$ can make $d=16^t$
extremely large, and the local-code search is exponential in $d^2$.
The proof absorbs these costs because $B\le\log n$ and every fixed
power of $\log n$ is eventually smaller than
$n/\sqrt{\log n}$. That comparison does not predict a usable
crossover at realistic input sizes. A practical implementation would
need substantially better local constructions and a careful
implementation of the source engine.

\section{Further research questions}
\label{sec:future}

\paragraph{Can the incidence information be compressed further?}
Our modified checks have $S=(B+2)^{O(1)}$ and therefore cost
$O(n\log(B+2))$ wires under the source's subset encoding. A
stronger reduction would need to encode or sum the relevant
incidence selections using fewer wires, or exploit dependencies
among their test coefficients. Merely suppressing the sparsity
parameter in the notation would not accomplish this. The sharper
bound \eqref{eq:wire-entropy}, involving the actual total incidence
count $N$, suggests studying nonuniform tests and many repeated
scopes.

\paragraph{Can one improve the check-support tradeoff?}
The explicit degree bound \eqref{eq:d-bound} is intentionally loose.
Better expander families and explicit local codes can reduce its
power and avoid exhaustive local search. Such changes could improve
constants and feasibility. As long as support remains a fixed power
of $B$, its logarithm remains $\Theta(\log B)$, so this alone would
not improve the asymptotic exponent scale proved here.

\paragraph{Can one sum weighted modes collectively?}
Proposition~\ref{prop:obstruction} rules out a purely numerical
improvement from adding only a total-degree cutoff. It leaves open
structured generating-function methods, algebraic compression of
packet tensors, or exact summation procedures that avoid treating
every allowed mode tuple separately. An improvement here would
affect the fixed-density bound as well.

\paragraph{Can density and sparsification be optimized jointly by a new invariant?}
The objective \eqref{eq:sparsify-objective} retains the
$\log\log n$ loss. To beat the unrestricted source bound through
this route, one would need stronger information about the leaves
than their maximum occurrence bound, such as a provably smaller
incidence entropy aggregated over the entire disjoint tree. A
uniform worst-case improvement has not been established here.

\paragraph{Can the extension be mechanized economically?}
A formal proof could reuse the source's tensor theorem and parity
lift as interfaces, then formalize the explicit graph powering,
entropy and distance inequalities, shortening, greedy elimination,
and the uniform asymptotic comparison. The exact-integer formulation
\eqref{eq:integer-tests} removes real arithmetic from parameter
selection. This is a proposed formalization path, not a claim of a
completed formal proof.

\begin{thebibliography}{9}
\bibitem{OAI2026}
OpenAI, \emph{The Exponential Time Hypothesis is False},
October 17, 2026, revised October 25, 2026.
Source snapshot: \texttt{openai/math}, commit
\texttt{7ef5c0cc5ca4b0e072b7a050fca09ab24bc6fa3a}.
See the main theorem; ``Initial tests that exclude small violation
sets''; ``Choosing additional equations by counting common zeros'';
and the appendices ``A sharper bound for exact counting'' and
``Proof of the contraction theorem.''
\href{https://github.com/openai/math/tree/7ef5c0cc5ca4b0e072b7a050fca09ab24bc6fa3a/preprints/The-Exponential-Time-Hypothesis-is-False-October-25-2026}{Permanent repository link}.

\bibitem{GG1981}
O. Gabber and Z. Galil,
\emph{Explicit constructions of linear-sized superconcentrators},
Journal of Computer and System Sciences \textbf{22} (1981),
407--420.
\href{https://doi.org/10.1016/0022-0000(81)90040-4}{doi:10.1016/0022-0000(81)90040-4}.

\bibitem{MGGFormal}
Archive of Formal Proofs,
\emph{Expander Graphs}, theory \texttt{Expander\_Graphs\_MGG},
section ``Margulis Gabber Galil Construction.''
The graph definition records the eight affine port maps;
\texttt{mgg\_numerical\_radius} establishes the spectral bound.
Accessed October 26, 2026.
\href{https://isa-afp.org/browser_info/current/AFP/Expander_Graphs/Expander_Graphs_MGG.html}{Formal theory and proof}.
\end{thebibliography}
\end{document}
